\documentclass[12pt,letterpaper]{article}
\usepackage[english]{babel}
\usepackage[margin=1in,headheight=15pt,headsep=19pt]{geometry}
\usepackage{amsmath}
\usepackage{graphicx}
\usepackage{fancyhdr}
\usepackage[hidelinks]{hyperref}
\usepackage{float}
\usepackage{booktabs}

\renewcommand{\rmdefault}{ptm}
\pagestyle{fancy}
\fancyhf{}
\fancyhead[L]{ECON1750}
\fancyhead[C]{Final Project: PCA Investing}
\fancyfoot[C]{\thepage}

\newcommand{\draftfill}[1]{\textbf{[\emph{#1}]}}
\newcommand{\answer}{\par\medskip}

\begin{document}

\begin{enumerate}
    \item Basic explanation of PCA, and the math foundations showing it is the eigenvectors of the covariance matrix. Then translate to the factor model from paper
    \item Directly answer each of the questions? And also put figures. We also need to answer them in an interesting way rather than just literally answering what the text is, i.e. the proportion of variance explained over time changes, and we can make a plot showing that. Some kind of more in depth analysis per question. Showing how much the components change over time, and also does removing outliers affect the calculated components.
    \item Explaining the methodology of the strategy framework we developed.
    \item Describe each strategy, what the intuition is, and what the results are.
\end{enumerate}

Strategy testing:
\begin{enumerate}
    \item Regime robustness, how well does it perform in each year and is one mega year carrying it.
    \item Parameter sensitivity: if we vary the parameters, then does the strategy remain profitable.
    \item Trading costs, showing how much it trades vs. the returns and how much of the edge is taken up by trading costs.
    \item Calibration: fitting a regression on the signal and the target is the actual observed return over the next holding period. A well-calibrated strategy shows that this is actually an increasing relationship.
    \item Holdout: test holdout, because it is better than claiming we got 400\% returns. Spencer will clown us if we do this.
    \item To answer the last question, sweep the strategies with different K values.
\end{enumerate}

\begin{itemize}
    \item Jerry: write about the strategy framework (answers how we set up question 7, and also question 11), find the holdout data
    \item Kayla: strategy testing on pairs trading and Treynor-Black
    \item Daniel: residual mean reversion strategy 
    \item Sameer: basic explanation of PCA, and the math foundations showing it is the eigenvectors of the covariance matrix, write up 1-6 in more interesting ways, plots, etc...
    \item Owen: math
\end{itemize}
\newpage

\section{PCA, explained}
Principal component analysis is a method of dimensionality reduction. For this project, we consider daily trading data for 1201 stocks over 24 years. Many stocks' returns are correlated, and we want to tease out the fundamentals, i.e. principal components, of these correlations. By finding the principal components, we can split out stock returns into two essential parts: a "shared part" representing how a stock reacted to common factors, and a "factor neutral" part representing everything else, which we can trade on. 
\\
\\
The first step is to standardize our returns to the form 
$$z_{i,t}=\frac{r_{i,t}-\bar{r}_i}{\bar{\sigma_i}}$$
such that returns have $E[Z_i]=0$ and $Var[z_i]=1$.



\newpage
In this project, we will have you build a quantitative trading strategy based
on standard machine-learning techniques. The project consists of two parts.
First, you will build (using PCA) a very simple model of stock alphas.
Second, you will design and back-test a trading strategy based on your model.

\begin{enumerate}
\item Using stock returns at the daily (or monthly) level, perform a principal
component analysis (PCA) on daily (monthly) stock returns. Feel free to refer
to online resources. How many factors are needed to explain a reasonable
proportion of the total variance in stock returns? (There is no right answer
to this, although there are standard share of variance measures that are
some diagnostic tests). (\emph{Hint:} what preprocessing should you do
before performing a PCA, and why?)

\answer
\answer
We performed principal component analysis (PCA) on daily returns for
1,201 stocks over 6,037 trading days from January 3, 2000, through
December 29, 2023. Before fitting PCA, we subtracted each stock's sample
mean and divided by its sample standard deviation. This gives each
stock unit variance and is equivalent to performing PCA on the return
correlation matrix.

Figure~\ref{fig:q1-explained-variance} shows a sharp drop after the
first principal component. PC1 explains 26.70\% of standardized return
variance, and the first five components explain \textbf{34.63\%}.
Components two through five add 7.93 percentage points combined.
We use five factors as a compact summary of the largest patterns in
returns. If the goal is to explain at least half of standardized return
variance, we need \textbf{70 components}. The number rises to 278 for
70\% and 672 for 90\%.

To see how preprocessing affects the leading component, we repeated
PCA on centered raw returns using the same stocks and dates. The 240
most volatile stocks accounted for 28.04\% of PC1's squared loadings
before standardization and 10.67\% afterward. If squared loadings were
distributed evenly across stocks, this group would account for 20.0\%.
Standardization therefore reduced the concentration of PC1's loadings
among the most volatile stocks in our sample. We use the standardized
fit for our main results to identify shared movements across stocks on
a common volatility scale.

\begin{figure}[H]
    \centering
    \includegraphics[width=0.85\linewidth]{figures/q1_explained_variance.png}
    \caption{\textbf{Explained variance from PCA on daily stock returns.}
    The left panel shows the share of variance explained by each of the
    first 25 principal components. The right panel shows cumulative
    explained variance, with the number of components on a logarithmic
    scale. The dashed vertical line marks \(K=5\), which explains
    34.63\% of the variance in standardized returns and 27.04\% in
    centered raw returns. Both fits use the same stocks and dates.}
    \label{fig:q1-explained-variance}
\end{figure}

\item Let's focus on a low number of factors: $K=5$. Compute (based on the total
sample for now) the $K=5$ principal components.

\begin{itemize}
\item Each principal component should be a vector of length $N$, the number
of stocks in your universe. How can one translate a principal component vector
$v_k$ into a tradable portfolio? (\emph{Hint:} what should you do with the
vector elements?)
\end{itemize}

\answer
We retain the five eigenvectors associated with the largest eigenvalues
of the full-sample return correlation matrix. Each eigenvector \(v_k\)
has 1,201 entries, one per stock. Together, the five components explain
34.63\% of standardized return variance.

To turn each component into an eigenportfolio, we divide each
coefficient \(v_k[i]\) by the stock's sample return standard deviation
\(\bar{\sigma}_i\). We then scale the resulting weights to one dollar
of gross exposure per dollar of capital:
\[
w_{k,i}
=
\frac{v_k[i]/\bar{\sigma}_i}
     {\displaystyle\sum_{j=1}^{1{,}201}
      \left|v_k[j]/\bar{\sigma}_j\right|}.
\]
Positive weights represent long positions, and negative weights represent
short positions. We orient each component so that its largest coefficient
in absolute value is positive. This fixes the sign convention for the
long and short positions shown in Figure~\ref{fig:q2-eigenportfolios}.

PC1 explains 26.70\% of standardized return variance and produces an
effectively long-only portfolio. Its ten largest positions account for
just 2.57\% of gross exposure, the lowest share among the five portfolios.
PC2 explains only 2.88\% of standardized variance, yet its ten largest
positions account for 9.32\% of gross exposure, the highest share of
the five. Under our sign convention, PC2 has a net long exposure of
51.3\%, while PC3 and PC5 have smaller net short exposures of 7.0\%
and 8.8\%, respectively.

The inverse-volatility adjustment also changes which stocks receive
the largest positions. For PC2 and PC5, only two of the ten stocks with
the largest absolute PCA coefficients remain in the top ten after the
adjustment. For PC1 and PC4, seven of ten remain. A stock can therefore
be prominent in the standardized component and receive a smaller
position in the eigenportfolio because of its return volatility.

\begin{figure}[H]
    \centering
    \includegraphics[width=0.85\linewidth]{figures/q2_eigenportfolios.png}
    \caption{\textbf{Largest positions in the five PCA eigenportfolios.}
    Each panel shows the six largest positions by absolute weight on a
    common scale. Positive and negative bars indicate long and short
    positions. Numeric labels are stock IDs. Panel titles report each
    component's share of standardized return variance.}
    \label{fig:q2-eigenportfolios}
\end{figure}


\item For each principal component vector $v_k$, compute the factor returns
for the eigenportfolio
$f_t^k=\sum_i\frac{v_k[i]}{\bar\sigma_i}\cdot r_{i,t}$, where
$\bar\sigma_i$ is the standard deviation for stock $i$ returns. Why do we
want to normalize the $v_k$ by the standard deviation?

\answer
Using the five full-sample eigenvectors from Q2, we compute daily factor
returns for all 6,037 trading days:
\[
f_t^k =
\sum_{i=1}^{1{,}201}
\frac{v_k[i]}{\bar{\sigma}_i}\,r_{i,t},
\qquad k=1,\ldots,5.
\]
We use the scaling specified in Q3. Applying Q2's gross-exposure
normalization would multiply each factor series by a constant.

We divide each coefficient by the stock's standard deviation because
the eigenvectors were estimated on standardized returns. For example,
if two stocks have the same PCA coefficient and one is twice as volatile,
the more volatile stock receives half the raw-return weight. A
one-standard-deviation move in either stock then makes the same
contribution to the factor.

The scaling also preserves the original PCA scores up to a constant.
If \(\bar r_i\) is stock \(i\)'s sample mean, then
\[
f_t^k-\bar f^k =
\sum_{i=1}^{1{,}201}
v_k[i]\frac{r_{i,t}-\bar r_i}{\bar{\sigma}_i}.
\]
The right-hand side is the PCA score, so the factor returns are
uncorrelated over the full estimation sample. In our data, applying the
eigenvector coefficients directly to raw returns gives a largest
absolute pairwise correlation of \textbf{0.5672}. Dividing by each
stock's standard deviation brings that value to \textbf{zero up to
numerical precision}.

Figure~\ref{fig:q3-factor-volatility} shows that factor volatility
varies even with fixed weights. PC1's 63-day volatility reached
67.25 factor units in December 2008, about 3.8 times its full-sample
daily standard deviation of 17.91. PC2 and PC5 also peaked in December
2008, while PC4 peaked in May 2020. These values are in factor units
under the scaling above.

\begin{figure}[H]
    \centering
    \includegraphics[width=0.85\linewidth]{figures/q3_factor_volatility.png}
    \caption{\textbf{Rolling volatility of the five PCA factor returns.}
    Each line shows the standard deviation of daily factor returns
    over the 63 trading days ending on the plotted date. The
    full-sample factor weights are fixed throughout.}
    \label{fig:q3-factor-volatility}
\end{figure}


\item Focus on the first factor return: what's the correlation of that with
respect to market returns?

\answer
We compare the first PCA factor with the daily value-weighted U.S.
market return from the
\href{https://mba.tuck.dartmouth.edu/pages/faculty/ken.french/data_library.html}
{Kenneth French Data Library}. We add the reported risk-free rate to
the market excess return to recover the total market return. The two
series match on all 6,037 trading days from January 3, 2000, through
December 29, 2023.

Under the sign convention used in Q2, the full-sample Pearson
correlation is \textbf{\(r=0.9373\)}. Regressing PC1 on the market
return with an intercept gives an \(R^2\) of 0.8785, so the market
return accounts for 87.85\% of PC1's daily variation in this regression.
This strong positive relationship is consistent with Q2's finding that
PC1 is an effectively long-only portfolio with weights spread across
many stocks.

Figure~\ref{fig:q4-pc1-market} shows that the strength of this
relationship varies over time. The trailing 252-trading-day correlation
has a median of 0.9623 and ranges from 0.7310 to 0.9898. The minimum
occurs in the window ending July 7, 2021, when PC1 tracked the broad
market substantially less closely than it did over the full sample.

\begin{figure}[H]
    \centering
    \includegraphics[width=0.85\linewidth]{figures/q4_pc1_market.png}
    \caption{\textbf{PC1 and the U.S. market return.}
    The left panel shows matched daily returns and the fitted
    regression line. The right panel shows trailing 252-trading-day
    Pearson correlations. The dashed line marks the full-sample
    correlation.}
    \label{fig:q4-pc1-market}
\end{figure}

\item For each stock \(i\), run the regression
\(r_{i,t}=\alpha_i+\sum_k\beta_{k,i}f_t^k+\epsilon_{i,t}\).
What is the interpretation of \(\alpha_i\)?

\answer
We regress each of the 1,201 stocks' daily returns on the five
PCA factor returns from Q3 over 6,037 trading days. The intercept
\(\alpha_i\) is the stock's fitted daily return when all five factor
returns are zero. Because the regression includes an intercept,
\[
\hat{\alpha}_i
=
\bar r_i-\sum_{k=1}^{5}\hat{\beta}_{k,i}\bar f^k.
\]
Thus, a positive alpha means the stock earned more on average than
the model attributes to its factor exposures. The residual
\(\hat{\epsilon}_{i,t}\) is the difference between the stock's actual
and fitted return on a particular day.

The median estimated alpha is \(0.298\) basis points per day, with
5th and 95th percentiles of \(-3.303\) and \(6.932\), respectively.
The median stock-level in-sample \(R^2\) is 0.3536.
In Figure~\ref{fig:q5-stock-alphas}, every stock lies below the
diagonal, meaning the fitted factor contributions were positive
on average for every stock. The median gap between mean return
and estimated alpha is \(5.846\) basis points per day.

The raw-return and alpha rankings have a Spearman correlation of
0.7264. Of the 240 stocks in the top quintile by raw average return,
\textbf{162 remain in the top quintile by alpha}. The factor
adjustment therefore changes 78 of the 240 names that a
top-quintile strategy would select. These estimates describe the
full sample. The out-of-sample backtest tests whether the alpha
ranking predicts future returns.

\begin{figure}[H]
    \centering
    \includegraphics[width=0.85\linewidth]{figures/q5_stock_alphas.png}
    \caption{\textbf{Estimated daily stock alphas.}
    The left panel shows the distribution of estimated alphas.
    The right panel plots each stock's estimated alpha against its
    mean daily return. The dashed diagonal marks equal values.
    Both return measures are in basis points per day.}
    \label{fig:q5-stock-alphas}
\end{figure}


\item How does this approach differ from standard factors
(i.e.\ using market, value, momentum)?

\answer
Standard factors are defined by portfolio construction rules. The
market factor is the broad market return minus the risk-free rate.
Value (HML) measures the return spread between stocks with high and
low book-to-market ratios, while momentum (Mom) measures the spread
between stocks with high and low past returns. PCA instead estimates
its factors directly from stock returns. It selects orthogonal
directions that explain as much in-sample variation in standardized
returns as possible. The economic interpretation comes after estimation.

We match the five PCA factors to daily market excess, value, and
momentum returns from the
\href{https://mba.tuck.dartmouth.edu/pages/faculty/ken.french/data_library.html}
{Kenneth French Data Library} on all 6,037 trading days. We compute
pairwise Pearson correlations and regress each PCA factor jointly on
the three standard factors with an intercept. We use the market
excess return here and the total market return in Q4.

Figure~\ref{fig:q6-pca-vs-standard-factors} shows that PC1 has a
correlation of \(0.9372\) with the market excess return. The three
standard factors jointly explain about 93.4\% of PC1's daily
variation, roughly 5.6 percentage points more than the market
excess return alone. Under our sign convention, PC3 has a
\(-0.6022\) correlation with HML, the strongest absolute value
correlation among the five components.

PC2 explains 2.88\% of standardized stock return variance, the
second-largest share among the components. Yet market, value, and
momentum jointly explain only \textbf{0.69\% of PC2's daily
variation}. The corresponding figure for PC5 is 3.73\%.
These results show that some PCA factors closely track familiar
benchmarks, while others capture patterns in this stock universe
that the three standard factors explain very little of.

\begin{figure}[H]
    \centering
    \includegraphics[width=0.85\linewidth]
        {figures/q6_pca_vs_standard_factors.png}
    \caption{\textbf{PCA factors and standard U.S. equity factors.}
    The left panel shows daily Pearson correlations. The right panel
    reports the in-sample \(R^2\) from regressing each PCA factor
    jointly on market excess, value, and momentum returns.}
    \label{fig:q6-pca-vs-standard-factors}
\end{figure}

\item Construct a trading strategy based on your estimate of $\alpha_i$.
If you are totally out of ideas, one thing would be to sort the stocks
based on your estimates of $\alpha_i$, go long stocks in the top quintile
of alpha, go short stocks in the bottom quintile of alpha. A trading
strategy should consist of how much of each stock $i$ to hold (can be
negative) at any point in time $t$. Note that if you put weights $w_{i,t}$
on stock $i$ at time $t$, your returns $r_{\mathrm{port},t+1}$ is given by
$\sum_i w_{i,t}\cdot r_{i,t+1}$.

\begin{itemize}
    \item Feel free to be more entrepreneurial/creative in this section.
    \item Make sure your trading strategy is not based on future returns,
    e.g.\ you go long stocks that will move up tomorrow.
\end{itemize}

\item Visualize / report the cumulative returns of the trading strategy.
Let $r_{\mathrm{port},t}$ be the returns of your trading strategy for
time $t$.

\item Compute the Sharpe ratio of your trading strategy.

\item Compute the information ratio of your trading strategy:

\begin{itemize}
    \item Formally, run the regression:
    $r_{\mathrm{port},t}=\alpha+\beta r_{\mathrm{mkt},t}$.
    Your reference return (for the information ratio) should be
    $\beta\cdot r_{\mathrm{mkt},t}$.
    \item Is your strategy market neutral?
\end{itemize}

\item Note that you have estimated the PCA using the entire sample.
Thus, your estimates of your trading performance should be somewhat
positively biased. Now, repeat your exercise, where you simulate a
trading strategy across time, where you a) estimate your factors only
based on data in the past (so at 2015 January, you construct your
PCAs using data up to 2015 January), and then construct a trading
strategy, and see how this strategy performs. A simple starting point
is to estimate $\alpha_i$s using data prior to 2015, and assume
$\alpha_i$s stay fixed afterwards. Of course, you can also update your
$\alpha_i$ estimates every day/month when you roll forward. Be very
clear how you avoided the look-ahead bias.

\item Why stop at $K=5$? Across multiple values of $K$, re-do all of
the above exercises. Which $K$ seems to work best? Is there a
significant difference between in-sample assessments (one where you
estimate the factors based on all the data) and out-sample assessments
(one where you estimate the factors only based on data available at
the time and roll forward)? Discuss.
\end{enumerate}

\small
\noindent\textbf{Reference.}
Kelly, B., Pruitt, S., and Su, Y. (2018).
\emph{Characteristics Are Covariances: A Unified Model of Risk and Return.}
October 15 working paper. \href{https://ssrn.com/abstract=3032013}{SSRN 3032013}.

\newpage

\input{framework}

\section{Trading Strategies}

We build two strategies on top of the PCA factor model. The main strategy, a \emph{contrarian factor-neutral Treynor--Black portfolio}, trades directly on the estimated alphas $\alpha_i$. The second, \emph{PCA pairs trading}, uses the principal component loadings to find pairs of economically similar stocks and trades the convergence of their prices.

The data are daily returns for $N = 1{,}201$ stocks from January 2000 to December 2023. Both strategies are simulated out-of-sample: every estimate is recomputed on a trailing window as the simulation moves forward, and trading begins in January 2002 once 500 days of history are available.

\subsection{Rolling factor model}

\begin{itemize}
    \item \textbf{Principal components.} Every 5 trading days, we run PCA on the correlation matrix of the previous 500 days of returns. Standardizing returns first prevents high-volatility stocks from dominating the components. We keep the top $K = 5$ eigenvectors $v_1, \dots, v_5$.
    \item \textbf{Factor returns.} Each eigenvector defines an eigenportfolio with return $f^k_t = \sum_i \frac{v_k[i]}{\bar\sigma_i} r_{i,t}$.
    \item \textbf{Alphas and betas.} Each day, for each stock, we regress the previous 250 days of returns on the five factor returns,
    \[
        r_{i,t} = \alpha_i + \sum_{k=1}^{K} \beta_{k,i} f^k_t + \epsilon_{i,t},
    \]
    which gives the alpha $\alpha_i$, the factor betas $\beta_{k,i}$, and the residual variance $\sigma^2(\epsilon_i)$.
\end{itemize}

\subsection{Strategy 1: Contrarian factor-neutral Treynor--Black}

\paragraph{Scores.} Under the Treynor--Black model, the optimal active weight on a stock is proportional to its alpha divided by its idiosyncratic variance: a large alpha on a very noisy stock deserves less weight than a moderate alpha on a stable one. We use the negative of this ratio,
\[
    s_i = -\frac{\alpha_i}{\sigma^2(\epsilon_i)}.
\]
The sign is flipped because, out of sample, stocks with high estimated alpha go on to \emph{underperform} (Table~\ref{tab:development}). Past alpha in this universe reverses: it reflects temporary mispricing or estimation noise, not persistent outperformance. Section~\ref{sec:holdout} confirms that this reversal is already present using only 2002--2017 data. To keep a few extreme stocks from dominating, we convert the scores to cross-sectional z-scores and cap them at $\pm 3$.

\paragraph{Factor neutrality.} We project the scores onto the space orthogonal to a vector of ones and the $K$ beta vectors, so the portfolio satisfies
\[
    \sum_i w_i = 0 \quad \text{(dollar neutral)}, \qquad \sum_i w_i \beta_{k,i} = 0 \quad \text{for } k = 1, \dots, K \quad \text{(factor neutral)}.
\]
Because the first principal component is essentially the market, this also makes the strategy approximately market neutral.

\paragraph{Scaling and trading.} Weights are scaled so gross exposure is $\sum_i |w_i| = 1$ (50\% long, 50\% short), and the portfolio is rebalanced daily.

\begin{table}[H]
\centering
\caption{Development of the strategy (out-of-sample, 2002--2023, 5bp fees unless stated).}
\label{tab:development}
\begin{tabular}{lrrrr}
\toprule
Strategy & Sharpe & Sharpe (no fees) & Return & Market beta \\
\midrule
Long top 20\% $\alpha$, short bottom 20\% & $-1.31$ & $-0.98$ & $-6.1\%$ & $-0.07$ \\
Contrarian 20\% sort & 0.66 & 0.98 & 3.1\% & 0.07 \\
Contrarian 20\% sort, factor-hedged & 0.72 & 1.12 & 2.9\% & 0.04 \\
Contrarian Treynor--Black, not hedged & 0.76 & 1.18 & 2.4\% & 0.07 \\
\textbf{Contrarian Treynor--Black, factor-hedged} & \textbf{0.85} & \textbf{1.31} & 2.5\% & 0.04 \\
\bottomrule
\end{tabular}
\end{table}

Each design choice helps. Reversing the sign turns a losing strategy into a profitable one. Weighting by $\alpha_i / \sigma^2(\epsilon_i)$ improves on equal weights. The factor hedge leaves the return almost unchanged (2.5\% vs.\ 2.4\%) but lowers volatility from 3.2\% to 2.9\% and market beta from 0.07 to 0.04: it removes unrewarded factor risk without sacrificing alpha. Market beta is measured against an equal-weighted average of all stocks in the universe.

\begin{figure}[H]
\centering
\includegraphics[width=\textwidth]{figures/cumulative_returns_tb.pdf}
\caption{Growth of \$1 invested in Contrarian Treynor--Black, with and without 5bp transaction costs. Shaded regions mark the 2008--09 financial crisis and the 2020 COVID crash.}
\label{fig:cumulative-tb}
\end{figure}



\subsection{Strategy 2: PCA pairs trading}

\paragraph{Pair selection.} Every 21 trading days, we describe each stock by its loadings on the $K$ principal components, $\ell_i = (v_1[i]\sqrt{\lambda_1}, \dots, v_K[i]\sqrt{\lambda_K})$. We then greedily match the closest stocks in this space, keeping the 50 best non-overlapping pairs. Because the two stocks in a pair have nearly identical factor exposure, the spread between them is approximately factor neutral and driven mostly by stock-specific news.

\paragraph{Trading rule.} For each pair $(i, j)$, we track the log price ratio $\log(P_i / P_j)$ over the previous 60 days and compute its z-score $z$. When $z < -2$ we buy $i$ and short $j$; when $z > 2$ we do the reverse. We close the position once $|z| < 0.5$. Each pair is allocated a fixed $1/50$ of capital, so idle pairs hold cash.

Even though no industry data is used, the selected pairs are economically sensible: utilities (WEC/DUK, ATO/NI), natural gas producers (EQT/CNX), tower REITs (AMT/SBAC) and apartment REITs (EQR/MAA). The principal components recover industry structure from returns alone.

\begin{figure}[H]
\centering
\includegraphics[width=\textwidth]{figures/cumulative_returns_pairs.pdf}
\caption{Growth of \$1 invested in PCA pairs trading, with and without 5bp transaction costs. Shaded regions mark the 2008--09 financial crisis and the 2020 COVID crash.}
\label{fig:cumulative-pairs}
\end{figure}

\subsection{Residual mean reversion strategy}


The underlying strategy for the residual mean reversion strategy has to do with the underlying belief that when stocks move down on d more than what we had predicted from the data from days 1,..., d-1, then we should buy these particular stocks as we believe that the market has undervalues/overreacted to these stocks (and that they have higher relative expected value to assets that have moved as expected in our regression or assets that have moved higher than our prediction from the regression. The difference between this strategy and the Treynor black method is the fact that Treynor Black relies on inspection on the alpha parameter while the residual snapback strategy relies on buying stocks that we believe are relatively undervalued due to an overreaction in the markets.

For Treynor–Black, we have used a 250-day regression window since the estimated alpha from this data directly determines the trading signal and a longer window of data would gives us many more observations and would reduce the influence of a  anomalies in the data. 

However, for residual snapback, the regression serves a different purpose. This regression is supposed to give us a baseline/prior against we need to measure these recent unexpected returns. A shorter window such as a 100-day window lets that baseline adapt more quickly to recent changes in the stock’s factor exposures, while also at the cost of estimates with higher standard devations.  However, due to the short term edge We therefore use 100 days for snapback as a design choice suited to its short-horizon signal.
\quad

Now, to understand the "signal" of the residual between the actual daily return compared to the regressions predictions, we need to get an understanding of how much of a deviation/size of residual is normal for a particular stock. For example, if a particular stock barely deviates from the regression prediction and then randomly is 5 percent different than what the regression predicts would have a different signal than a stock that on average has a 5 percent deviation from the regression prediction. To do this, we look at the last 60 trading days and fit the regression model to 1 day less than our current day. Then, we calculate the residuals on those 60 days. After we have these residuals, we use these 60 data points to get a sample variance for the residual as well as a sample mean. Then, we examine how the next three days of trading looks like after we have out 100 days of data. After normalizing the average residual by transforming it into a z score, we get a heuristic/z-score statistic for how much the stock has gone down compared to its "fair value." Although "fair value" is not the correct depicition since we have no way of understanding what is a "fair value." 

\[
\mathrm{Score}_{i,t}
=
\frac{\displaystyle\sum_{k=0}^{2}\varepsilon_{i,t-k}}
     {\sqrt{3}\,\widehat{\sigma}_{i,t}},
\qquad
\widehat{\sigma}_{i,t}
=
\sqrt{\frac{1}{59}
\sum_{k=3}^{62}
\left(\varepsilon_{i,t-k}-\overline{\varepsilon}_{i,t}\right)^2},
\qquad
\overline{\varepsilon}_{i,t}
=
\frac{1}{60}\sum_{k=3}^{62}\varepsilon_{i,t-k}.
\]

Note that the numerator here in the score function is the sum of the three errors (the days 1 day 2 day and 3 day after the 100 previous days of regression and the 60 previous days of understanding how our model is calibrated.) A negative value in the error refers to the stock underperformin g the return compared to what our factor model was predicting over those three days. On the other hand, a positive value means that we thought that the market price has "overshot" what the factor model would have predicted. If we assume that the expected value of each of these errors is 0 (the factor model is an unbiased estimator of what the price of a stock should be condicitioned upon the rest of the market prices), then we essentially know that the sum in the numerator is how much the price and the "condiitonal mean disagree." Now, the denominator of the stock is used to "standardize" the error, as we have mentioned earlier. We have estimated summa i,t hat from the previous 60 days of trading data and then we know that the sum of three random variables has a variance that is 3 times a single one, since the standard deviation is sqrt of variance, then we know that the sum of three std devs is equivalent in expansion to sqrt(3) * standard dev which is why we divide by sqrt(3) in the denominator.
(An assumption that we are making when dividing by sqrt(3) is that the residuals are uncorrelated, which may be a strech given that this "error" is hypothesized to be due to irrational market sentiment, which would likely presist and be correlated throughout the three days. If we were to think about this strategy more, one way to add complexity would be to alter the denominator to match this. However, since this score applied to all of the stocks in our portfolio, it is likely not that important if our signal is supposed to be multiplied by a particular magnitude that is constant/not dependent on our stock (given that a signal is only important as it has relative information, meaning that its used to compare two stocks). 

The second expression is how we have calulated the standard deviation. Hopefully, the average error is ei,t which would mean that we predicted the error well and then we are only dividing by 59 instead of 60 because the expression is sqrt(1/(n-1) * rss) for the sample standard deviation. 

Below here, I have also attached the 60-day residual mean across stocks, so that we can examine if it is true that that the regression model likely is an unbiased conditional estimate of a securities market price:  

\begin{figure}[H]
    \centering
    \includegraphics[width=0.5\linewidth]{errors.png}
    \caption{Factor-model residual diagnostic.}
    \label{fig:snapback-residual-errors}
\end{figure}

\begin{figure}[H]
    \centering
    \includegraphics[width=0.5\linewidth]{erorrs_mean.png}
    \caption{Rolling residual mean diagnostic.}
    \label{fig:snapback-rolling-errors}
\end{figure}

\begin{figure}[H]
    \centering
    \includegraphics[width=0.5\linewidth]{error_mean_.png}
    \caption{Additional residual mean diagnostic.}
    \label{fig:snapback-mean-error}
\end{figure}

Notice that after examining these figures, they show that the factor model has quite a small average signed error (0) across the stocks. 
 This suggests that when using the factor exposures and the regression used/fitted until the prior day, the model explains returns without a  tendency to overstate or understate them or understate them on average. 

 Trading Sub-Strategies:    

\paragraph{Naive residual snapback.}

On the last trading day of each week, we rank stocks by their residual
scores and select up to 20 stocks with "scores" of $-2$ or lower. If there are more than 20 stocks that have scores of -2 or lower, then we select the ones that have the lowest scores. If we have n stocks, we would go long in each stock with a proprtion of 1/n for each stock. If there are zero stocks that qualify, then we would hold all of our portfolio in cash. We re-evaluate our portfolio each trading day of the week by looking at the past 3 days, Wednesday throguh firsthand and the recalibrate our portfolio at the end of the week. 

\begin{figure}[H]
    \centering
    \includegraphics[width=0.5\linewidth]{naive_snapback.png}
    \caption{Naive Snapback cumulative net P\&L versus the equal-weight market.}
    \label{fig:snapback-naive-pnl}
\end{figure}

\begin{figure}[H]
    \centering
    \includegraphics[width=0.5\linewidth]{stricter_cutoff.png}
    \caption{Stricter-cutoff Snapback cumulative net P\&L versus the equal-weight market.}
    \label{fig:snapback-strict-pnl}
\end{figure}

\paragraph{Experimental Factor Hedging on Top of Naive Residual Snapback}

Although the signal may be choosing stocks that are expected to be higher than the market, it can still have exposures to those 5 PCA factors/vectors. For example, our portfolio could be heavily positively exposed to the 2nd principal component, if we have an equal weight long portfolio that is sensitive to these movements, then that would not be ideal. To hedge out our portfolio to these exposures, we can naively start with our selected group of stocks and adjust their weights. 

The fully hedged strategy aims to mitigate this risk by collecting the top 40 eligible stocks that have the "lowest scores" (meaning that they have moved downwards more than what we expect) and set these as our "long candidate space." Then, we take up to 40 stocks with the most positive scores as short candidates. The longs express the fact that we believe some of the candidates have moved too far down and should be worth more (with relation to the rest of the market) and the shorts we have are the stocks that we believe have moved too far up in relation to the other stocks in our portfolio space. We then start with target weights that would allow us to provide 50 percent of our capital to the shorts and 50 percent to the longs. 

%look into the below well

We can denote $\mathbf{w}^{(0)}_t$ denote these signed target
weights (which we have naivelly set it to be .5 * 1/(number of long candidates) for each vector entry that corresponds to the stock that has a negative stock and .5 * 1/n(number of short candidates) for a stock that corresponds to a positive score. Now, also let $B_t$ be the matrix whose columns contain the five estimated PCA factor loadings for each selected stock.  Note that the dimension of B is 5 * (number of stocks that are both in the long or short candidate space). 

For any weight
vector $\mathbf{w}_t$, the portfolio's estimated factor exposures are
$B_t\mathbf{w}_t$. We choose weights as close as possible to the
initial targets while imposing factor neutrality and preserving the
direction of each trade:
\[
\begin{aligned}
\min_{\mathbf{w}_t}\quad &
    \frac{1}{2}\|\mathbf{w}_t-\mathbf{w}^{(0)}_t\|_2^2 \\
\text{subject to}\quad &
    B_t\mathbf{w}_t=\mathbf{0}, \qquad
    \sum_i |w_{i,t}|=1,\\
&   w_{i,t}\geq 0 \quad \text{for long candidates},\\
&   w_{i,t}\leq 0 \quad \text{for short candidates}.
\end{aligned}
\]

After performing this strategy and keeping all other frameworks from the previous strategy the same, we get the following results: 

\begin{figure}[H]
    \centering
    \includegraphics[width=0.5\linewidth]{fully_hedged.png}
    \caption{Fully hedged Snapback cumulative net P\&L versus the equal-weight market.}
    \label{fig:snapback-hedged-pnl}
\end{figure}

\paragraph{Fixed-weight hybrid residual snapback.}
Our naive intitial Snapback portfolio gives full weight to (at most) 20 long-only
rebound signals, while the fully hedged portfolio offsets exposure to the five PCA factors by constructing the weights to meet a Principal component neutral portfolio. Now, what if we were to blend both these portfolios to retain the naive strategies upside (since this portfolio is known to perform better in terms of returns compared to market) while reducing exposure/risk through the hedged portfolio. 

We test whether blending the two
portfolios can retain some of the naive strategy's upside while
reducing its exposure to common factor moves.

%look into the below more: 

Essentially, we are keeping each of the two strategies exactly the
same. All we are doing is constructing both portfolios separately
and then combining their stock weights:
\[
    \mathbf{w}_t^{\mathrm{hybrid}}
    = \lambda\,\mathbf{w}_t^{\mathrm{naive}}
    + (1-\lambda)\,\mathbf{w}_t^{\mathrm{hedged}}.
\]

Our parameter that we have the freedom to change is $\lambda$,
which determines the mix of the two portfolios in our hybrid model.
We test $\lambda = 0.30$, $0.50$, and $0.70$. Note that
$\mathbf{w}_t^{\mathrm{naive}}$ contains the weights we have determined
for each stock based on our naive criteria, and
$\mathbf{w}_t^{\mathrm{hedged}}$ contains the weights we have determined
for each stock based on our hedging criteria.

\begin{figure}[H]
    \centering
    \includegraphics[width=0.5\linewidth]{30percentnaive.png}
    \caption{Fixed blend of 30\% naive and 70\% hedged.}
    \label{fig:snapback-hybrid-30}
\end{figure}

\begin{figure}[H]
    \centering
    \includegraphics[width=0.5\linewidth]{50percentnaive_50hedged.png}
    \caption{Fixed blend of 50\% naive and 50\% hedged.}
    \label{fig:snapback-hybrid-50}
\end{figure}

\begin{figure}[H]
    \centering
    \includegraphics[width=0.5\linewidth]{70percentnaive_30percenthedged.png}
    \caption{Fixed blend of 70\% naive and 30\% hedged.}
    \label{fig:snapback-hybrid-70}
\end{figure}

\paragraph{Volatility-weighted hybrid residual snapback.}

Note that in the fixed-weight hybrid strategies, we are supposed to choose a value of $\lambda$
and keep this parameter the same throughout the trading horizon. Note that this may not be the ideal way to construct this strategy. For example, if we are in an uncertain time where stocks will tend to move around a lot, then it is better to have a portfolio that is more hedged to these risks. However, if we believe that stocks have lower volatility and that the tail risk of a particular asset class/sector/principle component dropping sharply in value, we are willing to take a more aggressive holding of stocks. 

We plan on changing this parameter depending on the volatility of the stocks. 

The first thing that we do is we calculate an equal-weight market return by averaging stock returns across our entire csv file for each day. (Essentially we calculate the return of the  equal weight market portfolio). Then, we can also estimate the sample standard deviation of these returns based off a specific set of returns. Thus, to estimate daily market volatility for each date $t$, we measure the sample standard deviation of these returns over the trailing 30 trading days. Note that this estimate uses only returns that are available by day $t$. 

After we do this, we let the variable $p_t \in [0,1]$ be the
resulting historical percentile rank of volatility. For example, if this is the lowest volatility in all of the history, then we would let $p_t$ be 0. However, if this was a the highest volatility, then we would let $p_t$ be 1. 

We will now use this $p_t$ value as the weight of our portfolio that we would like to have in the fully hedged portfolio. In other words, the weight on the naive portfolio is $\lambda_t = 1-p_t$. 

Notice that the following hedged allocation for the weights are: 

Although this graph looks pretty good, something to notice is that at the beggining of the timeline, we are alternative from 0 percent to 100 percent in the weights. This is because the rank is very sensitive to the few data points that we have. For example, at the beggining of this process we would have a sparse amount of data (for example the first data point is both the lowest and the highest volatility seen yet, so it is hard to factor this information into the model). 

To avoid this problem of letting a small number of early observations drive the allocation, we can use a symmetric Laplace smoothing to our weight allocation. 
Before each weekly re balance, we compare the current estimate of volatility with all earlier 30-day volatility estimates. 
We define the following variables: $n_t$
be the number of earlier readings and $k_t$ be the number
below the current reading, with ties counting as 1/2. 

Then, we will add 126 pseudo-readings on each side of our current estimate. This gives us the following weight on the hedged portfolio:
\[
    \widetilde p_t = \frac{k_t+126}{n_t+252}.
\]

Note that now, with no earlier readings, the rule starts at a 50/50 mix. The starting mix receives weight of $252/(n_t+252)$, so its "bias" decreases as more observations accumulate. For example, with 252 earlier readings, a raw volatility percentile of 90\% becomes a hedge weight of 70\%, halfway between the raw percentile and the 50\% starting weight.

The final portfolio allocates $\widetilde p_t$ to the fully hedged
portfolio and $1-\widetilde p_t$ to the naive portfolio:
\[
    \mathbf{w}_t^{\mathrm{dynamic}}
    =
    (1-\widetilde p_t)\mathbf{w}_t^{\mathrm{naive}}
    +
    \widetilde p_t\mathbf{w}_t^{\mathrm{hedged}}.
\]

\begin{figure}[H]
    \centering
    \includegraphics[width=0.5\linewidth]{volatility_allocation.png}
    \caption{Weekly hedge allocation under the original volatility rule.}
    \label{fig:snapback-raw-allocation}
\end{figure}

Below is an image depicting the new weight allocation graph:

\begin{figure}[H]
    \centering
    \includegraphics[width=0.5\linewidth]{updated_vol_weights.png}
    \caption{Trailing 30-day volatility and smoothed weekly hedge allocation.}
    \label{fig:snapback-smoothed-allocation}
\end{figure}

\begin{figure}[H]
    \centering
    \includegraphics[width=0.5\linewidth]{dynamic_volatility.png}
    \caption{Dynamic Volatility Snapback cumulative net P\&L versus the equal-weight market.}
    \label{fig:snapback-dynamic-pnl}
\end{figure}

\subsection{Backtest design}

\paragraph{Timing and look-ahead bias.} Weights chosen at the close of day $t$ use only returns up to and including day $t$, and they earn the returns of day $t+1$: $r_{\text{port},t+1} = \sum_i w_{i,t} r_{i,t+1}$. We verified this directly: re-running the backtest on data truncated at December 2011 leaves every earlier position unchanged, up to floating-point rounding of order $10^{-13}$.

\paragraph{Transaction costs.} The assignment defines returns without costs, so we report results both without fees and with a cost of 5 basis points per dollar traded.

\paragraph{Survivorship bias.} Every stock in the sample trades continuously from 2000 to 2023, so the universe consists of firms that survived the full period. Firms that failed or were acquired are missing. This likely flatters both strategies, since buying a stock after a fall never ends in a delisting in this sample. It is a limitation of the data, not of the strategy code.

\section{Strategy Testing}

We test both strategies along six dimensions. Unless stated otherwise, all results are out-of-sample, from January 2002 to December 2023.

\subsection{Regime robustness}

\begin{table}[H]
\centering
\caption{Annual returns and Sharpe ratios.}
\label{tab:yearly}
\begin{tabular}{lrrrrrr}
\toprule
& \multicolumn{3}{c}{Contrarian Treynor--Black} & \multicolumn{3}{c}{PCA pairs trading} \\
\cmidrule(lr){2-4} \cmidrule(lr){5-7}
Year & Return & Return (5bp) & Sharpe & Return & Return (5bp) & Sharpe \\
\midrule
2002 & 5.6\% & 4.2\% & 1.93 & 13.8\% & 13.0\% & 3.52 \\
2003 & 3.6\% & 2.2\% & 1.64 & 0.9\% & 0.1\% & 0.33 \\
2004 & 2.6\% & 1.2\% & 1.51 & 2.6\% & 1.9\% & 1.02 \\
2005 & 0.7\% & $-0.7\%$ & 0.40 & 1.4\% & 0.7\% & 0.69 \\
2006 & 2.4\% & 1.1\% & 1.29 & 1.6\% & 0.8\% & 0.85 \\
2007 & $-1.1\%$ & $-2.4\%$ & $-0.38$ & 4.8\% & 4.0\% & 1.77 \\
2008 & 15.6\% & 13.7\% & 3.00 & 7.2\% & 6.2\% & 1.53 \\
2009 & 23.4\% & 21.5\% & 4.83 & 5.5\% & 4.8\% & 2.02 \\
2010 & 2.2\% & 0.8\% & 1.14 & 4.8\% & 4.1\% & 2.49 \\
2011 & 2.3\% & 0.8\% & 0.90 & 2.1\% & 1.3\% & 1.12 \\
2012 & 5.3\% & 3.9\% & 2.58 & $-0.6\%$ & $-1.2\%$ & $-0.42$ \\
2013 & 0.2\% & $-1.1\%$ & 0.14 & 1.8\% & 1.1\% & 0.94 \\
2014 & 2.7\% & 1.3\% & 1.62 & 1.1\% & 0.4\% & 0.79 \\
2015 & $-0.3\%$ & $-1.6\%$ & $-0.14$ & 0.1\% & $-0.6\%$ & 0.07 \\
2016 & 5.5\% & 4.1\% & 2.50 & 7.8\% & 7.0\% & 3.80 \\
2017 & 0.3\% & $-0.9\%$ & 0.14 & 3.9\% & 3.2\% & 2.28 \\
2018 & 1.7\% & 0.4\% & 0.70 & 3.6\% & 2.8\% & 2.06 \\
2019 & 5.1\% & 3.7\% & 2.46 & 4.5\% & 3.9\% & 2.71 \\
2020 & 8.1\% & 6.5\% & 1.26 & 9.5\% & 8.6\% & 1.91 \\
2021 & 1.6\% & 0.1\% & 0.61 & 2.7\% & 2.1\% & 1.14 \\
2022 & 1.6\% & 0.2\% & 0.44 & 0.7\% & $-0.1\%$ & 0.32 \\
2023 & $-0.8\%$ & $-2.1\%$ & $-0.22$ & 2.2\% & 1.6\% & 1.05 \\
\bottomrule
\end{tabular}
\\[4pt]
\small Returns and Sharpe ratios are without fees unless marked 5bp.
\end{table}

\begin{table}[H]
\centering
\caption{Concentration and stability of returns.}
\label{tab:regime}
\begin{tabular}{lcc}
\toprule
Measure & Treynor--Black & Pairs trading \\
\midrule
Years positive (no fees / 5bp) & 19 / 16 of 22 & 21 / 19 of 22 \\
Share of total profit from best year & 25\% (2009) & 16\% (2002) \\
Share of total profit from best two years & 43\% (2008--09) & 28\% (2002, 2020) \\
Sharpe, all years (no fees / 5bp) & 1.31 / 0.85 & 1.45 / 1.17 \\
Sharpe excluding best two years (no fees / 5bp) & 0.91 / 0.41 & 1.28 / 0.97 \\
Sharpe 2002--2012 / 2013--2023 (5bp) & 1.41 / 0.32 & 1.18 / 1.16 \\
Maximum drawdown (no fees / 5bp) & $-5.2\%$ / $-6.1\%$ & $-3.6\%$ / $-3.9\%$ \\
\bottomrule
\end{tabular}
\end{table}

No single year carries either strategy, but the two differ markedly (Figures~\ref{fig:cumulative-tb}--\ref{fig:annual-pairs}). Treynor--Black earns its best returns in crises: in 2008 it gained 15.6\% while the market fell 27\%, followed by 23.4\% in 2009 and 8.1\% in 2020. These are periods of panic selling and overreaction, when betting on reversal should pay. However, 2008--09 account for 43\% of its total profit, and its Sharpe ratio after fees fell from 1.41 in 2002--2012 to 0.32 in 2013--2023. Pairs trading is far more stable: profitable in 21 of 22 years, no more than 28\% of profit from any two years, and an essentially identical Sharpe ratio in both halves of the sample.

\begin{figure}[H]
\centering
\includegraphics[width=\textwidth]{figures/annual_returns_tb.pdf}
\caption{Annual returns of Contrarian Treynor--Black (no fees).}
\label{fig:annual-tb}
\end{figure}

\begin{figure}[H]
\centering
\includegraphics[width=\textwidth]{figures/annual_returns_pairs.pdf}
\caption{Annual returns of PCA pairs trading (no fees).}
\label{fig:annual-pairs}
\end{figure}

\subsection{Parameter sensitivity}

We vary one parameter at a time around the defaults. Because a 1000-day PCA window needs four years of history, all variants are compared over the common period beginning December 2003. The last two columns split this period into a selection period (through 2017) and a holdout period (2018--2023); see Section~\ref{sec:holdout}.

\begin{table}[H]
\centering
\caption{Parameter sensitivity: contrarian Treynor--Black.}
\label{tab:sensitivity-tb}
\begin{tabular}{llrrrrr}
\toprule
Parameter & Value & Sharpe & Sharpe (5bp) & Return (5bp) & Selection & Holdout \\
\midrule
Default & --- & 1.27 & 0.81 & 2.4\% & 1.08 & 0.40 \\
Alpha window (days) & 60 & 1.93 & 1.00 & 2.8\% & 1.13 & 0.83 \\
& 125 & 1.55 & 0.93 & 2.7\% & 1.13 & 0.64 \\
& 500 & 1.15 & 0.79 & 2.2\% & 0.88 & 0.64 \\
PCA window (days) & 250 & 1.25 & 0.78 & 2.3\% & 1.06 & 0.37 \\
& 1000 & 1.17 & 0.72 & 2.1\% & 0.96 & 0.36 \\
Score cap & 2 & 1.27 & 0.81 & 2.4\% & 1.10 & 0.39 \\
& 5 & 1.27 & 0.80 & 2.3\% & 1.06 & 0.41 \\
\bottomrule
\end{tabular}
\\[4pt]
\small Defaults: alpha window 250 days, PCA window 500 days, score cap 3. Selection and holdout columns include 5bp fees.
\end{table}

\begin{table}[H]
\centering
\caption{Parameter sensitivity: PCA pairs trading.}
\label{tab:sensitivity-pairs}
\begin{tabular}{llrrrrr}
\toprule
Parameter & Value & Sharpe & Sharpe (5bp) & Return (5bp) & Selection & Holdout \\
\midrule
Default & --- & 1.36 & 1.07 & 2.6\% & 1.04 & 1.15 \\
Entry threshold $|z|$ & 1.5 & 1.59 & 1.21 & 3.4\% & 1.20 & 1.25 \\
& 2.5 & 1.21 & 1.01 & 2.0\% & 0.92 & 1.19 \\
& 3 & 1.15 & 1.03 & 1.5\% & 0.99 & 1.14 \\
Exit threshold $|z|$ & 0 & 1.25 & 0.98 & 2.5\% & 0.97 & 1.02 \\
& 1 & 1.49 & 1.15 & 2.6\% & 1.05 & 1.37 \\
Spread window (days) & 20 & 1.93 & 1.43 & 3.4\% & 1.63 & 1.05 \\
& 40 & 1.59 & 1.25 & 3.1\% & 1.27 & 1.20 \\
& 120 & 1.29 & 1.08 & 2.6\% & 1.06 & 1.14 \\
Number of pairs & 20 & 0.68 & 0.48 & 1.7\% & 0.43 & 0.58 \\
& 100 & 1.73 & 1.38 & 2.7\% & 1.51 & 1.20 \\
Re-selection interval (days) & 5 & 1.54 & 1.14 & 2.5\% & 1.06 & 1.32 \\
& 63 & 1.42 & 1.19 & 3.1\% & 1.03 & 1.55 \\
PCA window (days) & 250 & 1.21 & 0.88 & 2.0\% & 0.76 & 1.13 \\
& 1000 & 1.54 & 1.28 & 3.3\% & 1.10 & 1.66 \\
Selection method & correlation & 1.89 & 1.35 & 1.6\% & 1.37 & 1.33 \\
\bottomrule
\end{tabular}
\\[4pt]
\small Defaults: entry 2, exit 0.5, spread window 60 days, 50 pairs, re-selection every 21 days, PCA window 500 days, selection by PCA loadings.
\end{table}

Both strategies remain profitable after fees under every variant tested (Figures~\ref{fig:sensitivity-tb} and~\ref{fig:sensitivity-pairs}). Pairs trading is robust, with most Sharpe ratios between 0.9 and 1.4 after fees. The one weak variant is trading only 20 pairs (0.48), which concentrates risk in too few positions. Selecting pairs by return correlation gives a similar Sharpe ratio but a much lower return, partly because it selects two share classes of the same company (e.g.\ Brown-Forman, Berkshire Hathaway), which barely diverge. For Treynor--Black, shorter alpha windows perform better. This suggests that part of its edge comes from short-term reversal in the most recent returns entering the estimation window.

\begin{figure}[H]
\centering
\includegraphics[width=0.7\textwidth]{figures/parameter_sensitivity_tb.pdf}
\caption{Sharpe ratio (5bp fees) of Contrarian Treynor--Black for each parameter variant. The vertical line marks the default settings.}
\label{fig:sensitivity-tb}
\end{figure}

\begin{figure}[H]
\centering
\includegraphics[width=0.7\textwidth]{figures/parameter_sensitivity_pairs.pdf}
\caption{Sharpe ratio (5bp fees) of PCA pairs trading for each parameter variant. The vertical line marks the default settings.}
\label{fig:sensitivity-pairs}
\end{figure}

\subsection{Trading costs}

\begin{table}[H]
\centering
\caption{Turnover and the cost of trading.}
\label{tab:costs}
\begin{tabular}{lcc}
\toprule
& Treynor--Black & Pairs trading \\
\midrule
Daily turnover & 10.9\% & 5.6\% \\
Annual turnover & 28$\times$ & 14$\times$ \\
Gross return (annual mean) & 3.8\% & 3.6\% \\
Fee drag at 5bp (annual) & 1.4\% & 0.7\% \\
Share of edge lost at 5bp & 36\% & 20\% \\
Break-even fee & 14.0bp & 25.7bp \\
\midrule
Sharpe at 0bp & 1.31 & 1.45 \\
Sharpe at 2bp & 1.13 & 1.33 \\
Sharpe at 5bp & 0.85 & 1.17 \\
Sharpe at 10bp & 0.38 & 0.88 \\
Sharpe at 20bp & $-0.57$ & 0.32 \\
\bottomrule
\end{tabular}
\end{table}

Figures~\ref{fig:fees-tb} and~\ref{fig:fees-pairs} show how each strategy's Sharpe ratio falls with costs. Treynor--Black rebalances daily and turns over its portfolio 28 times a year, so costs take about a third of its edge, and it stops being profitable at around 14bp per trade. Pairs trading holds each position until the spread converges, halving turnover. It loses only a fifth of its edge to costs and survives fees up to about 26bp.

\begin{figure}[H]
\centering
\includegraphics[width=0.65\textwidth]{figures/fee_sensitivity_tb.pdf}
\caption{Sharpe ratio of Contrarian Treynor--Black as a function of the transaction cost per dollar traded.}
\label{fig:fees-tb}
\end{figure}

\begin{figure}[H]
\centering
\includegraphics[width=0.65\textwidth]{figures/fee_sensitivity_pairs.pdf}
\caption{Sharpe ratio of PCA pairs trading as a function of the transaction cost per dollar traded.}
\label{fig:fees-pairs}
\end{figure}

\subsection{Calibration}

To check that each signal predicts returns in the intended direction, we compare the signal on day $t$ with realized returns over the next $h$ days. For each day, we regress the realized return on the cross-sectionally standardized signal. Table~\ref{tab:calibration-slopes} reports the average slope with Newey--West $t$-statistics to account for overlapping horizons.

\begin{table}[H]
\centering
\caption{Treynor--Black: mean excess return by signal decile (bp).}
\label{tab:calibration-tb}
\begin{tabular}{lrrrrrrrrrr}
\toprule
Horizon & D1 & D2 & D3 & D4 & D5 & D6 & D7 & D8 & D9 & D10 \\
\midrule
1 day & $-3.0$ & $-2.0$ & $-1.3$ & $-1.0$ & 0.6 & 2.1 & 1.8 & 1.6 & 1.0 & 0.1 \\
5 days & $-10.7$ & $-6.6$ & $-1.7$ & $-1.0$ & 3.1 & 6.2 & 3.6 & 3.6 & 2.9 & 0.7 \\
20 days & $-34.4$ & $-17.4$ & $-0.5$ & 2.1 & 14.2 & 22.1 & 10.0 & 6.9 & 3.4 & $-6.1$ \\
\bottomrule
\end{tabular}
\\[4pt]
\small D1 is the decile the strategy shorts most (highest alpha), D10 the decile it buys most. Excess return is relative to the cross-sectional average.
\end{table}

\begin{table}[H]
\centering
\caption{Pairs trading: mean change in the pair's log price ratio by spread z-score (bp).}
\label{tab:calibration-pairs}
\begin{tabular}{lrrrr}
\toprule
Spread z-score & Next 1 day & Next 5 days & Next 20 days & Observations \\
\midrule
$z < -3$ & 48.4 & 108.4 & 157.9 & 2,251 \\
$-3$ to $-2$ & 17.7 & 20.8 & 43.9 & 16,811 \\
$-2$ to $-1$ & 7.0 & 13.5 & 8.3 & 56,696 \\
$-1$ to $0$ & 3.0 & 5.0 & $-0.8$ & 64,245 \\
$0$ to $1$ & $-3.5$ & $-6.3$ & $-31.8$ & 64,170 \\
$1$ to $2$ & $-7.2$ & $-14.3$ & $-53.1$ & 55,074 \\
$2$ to $3$ & $-18.3$ & $-46.2$ & $-101.7$ & 15,567 \\
$z > 3$ & $-59.7$ & $-143.8$ & $-267.5$ & 2,036 \\
\bottomrule
\end{tabular}
\end{table}

\begin{table}[H]
\centering
\caption{Return per one standard deviation of signal.}
\label{tab:calibration-slopes}
\begin{tabular}{lrrrr}
\toprule
& \multicolumn{2}{c}{Treynor--Black} & \multicolumn{2}{c}{Pairs trading} \\
\cmidrule(lr){2-3} \cmidrule(lr){4-5}
Horizon & Slope & $t$-stat & Slope & $t$-stat \\
\midrule
1 day & 1.1bp & 5.2 & 9.6bp & 13.3 \\
5 days & 3.4bp & 3.4 & 19.3bp & 7.3 \\
20 days & 7.7bp & 1.8 & 39.0bp & 4.7 \\
\bottomrule
\end{tabular}
\end{table}

The pairs signal is very well calibrated (Figures~\ref{fig:calibration-tb} and~\ref{fig:calibration-pairs}). Future spread returns fall steadily with the z-score in every bucket and at every horizon, and the relationship is highly significant ($t = 13.3$ at one day). The Treynor--Black signal is significant at short horizons ($t = 5.2$) but only partly monotonic. Its predictive power is concentrated in the stocks it shorts: the highest-alpha decile underperforms strongly and persistently. The lowest-alpha stocks it buys most do not outperform, and at 20 days they underperform. In other words, high past alpha reverses, but very low past alpha does not. The edge also fades with horizon, consistent with the strategy's need for daily rebalancing.

\begin{figure}[H]
\centering
\includegraphics[width=0.65\textwidth]{figures/calibration_tb.pdf}
\caption{Calibration of Contrarian Treynor--Black: next 5-day excess return by signal decile.}
\label{fig:calibration-tb}
\end{figure}

\begin{figure}[H]
\centering
\includegraphics[width=0.65\textwidth]{figures/calibration_pairs.pdf}
\caption{Calibration of PCA pairs trading: next-day change in the pair log price ratio by spread z-score.}
\label{fig:calibration-pairs}
\end{figure}

\subsection{Holdout}
\label{sec:holdout}

To check that our design choices are not overfit, we treat 2018--2023 as a holdout. Any choice is made using only 2002--2017 data and then evaluated once on the holdout.

\begin{table}[H]
\centering
\caption{Selection period vs.\ holdout (Sharpe ratios, 5bp fees).}
\label{tab:holdout}
\begin{tabular}{llrr}
\toprule
Strategy & Choice & 2002--2017 & 2018--2023 \\
\midrule
Treynor--Black & Default settings & 1.08 & 0.40 \\
& Best setting on 2002--2017 (alpha window 60) & 1.13 & 0.83 \\
& Best $K$ on 2002--2017 ($K = 20$) & 1.11 & 0.63 \\
Pairs trading & Default settings & 1.04 & 1.15 \\
& Best setting on 2002--2017 (spread window 20) & 1.63 & 1.05 \\
& Best $K$ on 2002--2017 ($K = 15$) & 1.41 & 0.58 \\
\midrule
Following alpha (no fees) & Sign check & $-1.63$ & $-0.77$ \\
\bottomrule
\end{tabular}
\end{table}

Three results stand out. First, the contrarian direction is justified using the selection period alone: following alpha loses money in 2002--2017 (Sharpe $-1.63$), so reversing the sign was not chosen by looking at the full sample. Second, pairs trading holds up in the holdout (Sharpe 1.15, versus 1.04 before). Optimizing its parameters on the selection period did \emph{not} help: every optimized choice did worse in the holdout than the defaults, the classic signature of overfitting. Third, Treynor--Black weakens substantially in the holdout (Sharpe 0.40 with default settings), consistent with the decline in its second-half performance in Table~\ref{tab:regime}.

\subsection{Number of factors $K$}

Finally, we re-run both strategies for $K \in \{1, 2, 3, 5, 8, 10, 15, 20\}$, both out-of-sample and in-sample. In the in-sample version, the PCA and the alphas are estimated once on the full 2000--2023 sample and then traded over the same period, which introduces look-ahead bias.

\begin{table}[H]
\centering
\caption{Sharpe ratios by number of factors $K$.}
\label{tab:k}
\begin{tabular}{rrrrrrrr}
\toprule
& \multicolumn{4}{c}{Treynor--Black} & \multicolumn{3}{c}{Pairs trading} \\
\cmidrule(lr){2-5} \cmidrule(lr){6-8}
& \multicolumn{2}{c}{Out-of-sample} & \multicolumn{2}{c}{In-sample} & \multicolumn{2}{c}{Out-of-sample} & In-sample \\
\cmidrule(lr){2-3} \cmidrule(lr){4-5} \cmidrule(lr){6-7}
$K$ & No fees & 5bp & Contrarian & Follow $\alpha$ & No fees & 5bp & No fees \\
\midrule
1 & 0.95 & 0.68 & $-1.54$ & 1.54 & 0.78 & 0.52 & 0.62 \\
2 & 0.86 & 0.54 & $-2.12$ & 2.12 & 1.39 & 1.11 & 0.69 \\
3 & 0.87 & 0.51 & $-2.13$ & 2.13 & 1.13 & 0.86 & 1.13 \\
5 & 1.31 & 0.85 & $-2.18$ & 2.18 & 1.45 & 1.17 & 1.93 \\
8 & 1.46 & 0.92 & $-2.76$ & 2.76 & 1.41 & 1.09 & 2.18 \\
10 & 1.49 & 0.91 & $-2.77$ & 2.77 & 1.53 & 1.18 & 2.40 \\
15 & 1.58 & 0.94 & $-3.29$ & 3.29 & 1.51 & 1.14 & 2.39 \\
20 & 1.66 & 0.96 & $-3.42$ & 3.42 & 1.60 & 1.18 & 2.15 \\
\bottomrule
\end{tabular}
\end{table}

\begin{figure}[H]
\centering
\includegraphics[width=0.65\textwidth]{figures/k_sweep_tb.pdf}
\caption{Sharpe ratio (no fees) of Contrarian Treynor--Black by number of factors $K$, out-of-sample vs.\ in-sample. The in-sample version follows alpha rather than trading against it.}
\label{fig:k-tb}
\end{figure}

\begin{figure}[H]
\centering
\includegraphics[width=0.65\textwidth]{figures/k_sweep_pairs.pdf}
\caption{Sharpe ratio (no fees) of PCA pairs trading by number of factors $K$, out-of-sample vs.\ in-sample.}
\label{fig:k-pairs}
\end{figure}

\paragraph{In-sample vs.\ out-of-sample.} The difference is dramatic (Figures~\ref{fig:k-tb} and~\ref{fig:k-pairs}). In-sample, following alpha appears highly profitable, with Sharpe ratios of 2 to 3.4 that rise as more factors are added. This is look-ahead bias: the full-sample alpha is estimated from the very returns being traded, so stocks with high alpha mechanically earned high returns over the period. Out-of-sample, the same signal points the \emph{opposite} way: high past alpha predicts low future returns. An in-sample assessment would not just overstate performance; it would recommend the wrong strategy. Pairs trading is also inflated in-sample (Sharpe up to 2.4 vs.\ about 1.5 out-of-sample), though its direction is unchanged.

\paragraph{Which $K$ works best?} Out-of-sample, performance improves up to roughly $K = 5$--$8$ and then levels off, especially after fees. For Treynor--Black, the Sharpe ratio after fees is 0.85 at $K = 5$ and only 0.96 at $K = 20$. For pairs trading, every $K \ge 2$ gives a Sharpe ratio between 0.86 and 1.18 after fees. Picking $K$ by its past performance also did not carry over to the holdout (Table~\ref{tab:holdout}). We therefore keep $K = 5$: it captures most of the benefit, and larger values add estimation noise without a reliable improvement.

\section{Holdout results}
\label{sec:final-holdout}

The holdout covers 723 trading days from January 2, 2024 to November 17,
2026. Over this period, factor-hedged Treynor--Black returns -4.13\%
after fees with a Sharpe ratio of -0.53, while PCA pairs earns
4.03\% with a Sharpe ratio of 0.81. The difference comes from
both trading performance and costs: Treynor--Black generates little profit
before fees and trades frequently, whereas pairs earns more with lower
turnover.

We carry forward the window lengths selected on 2002--2017 and checked on
2018--2023, retaining $K=5$ from the earlier analysis. Each strategy runs with
one fixed configuration:
\begin{itemize}
    \item \textbf{Contrarian, factor-hedged Treynor--Black.} A 60-day
    factor regression, daily refits, and a score cap of three standard
    deviations. Targets are dollar neutral and factor neutral, with unit
    gross exposure.
    \item \textbf{PCA pairs.} A 20-day spread window, 50 disjoint
    loading-based pairs, reselection every 21 days using a 250-day formation
    window, and entry/exit thresholds of 2 and 0.5. Idle allocations remain
    in cash.
\end{itemize}
Both use a 500-day PCA window, five-day PCA refits, and fees of 5 basis
points per dollar traded. Portfolios start in cash, with models initialized
from the final 500 observations of the original sample. As in the earlier
tests, signals include the current day's return and trades execute at that
close.

\subsection{Performance and interpretation}

Table~\ref{tab:final-holdout} reports net performance under the framework's
gap-aware accounting. Pairs achieves a positive Sharpe ratio and a smaller
drawdown, although its annualized return remains modest.

\begin{table}[htbp]
\centering
\small
\begin{tabular}{lrrrr}
\toprule
Strategy & Net return & CAGR & Sharpe & Max DD \\
\midrule
Hedged TB, 60d & -4.13\% & -1.46\% & -0.53 & -4.91\% \\
PCA pairs, 20d & 4.03\% & 1.39\% & 0.81 & -1.75\% \\
\bottomrule
\end{tabular}
\caption{Holdout performance of the selected strategies, with $K=5$ and
5 bp trading fees. CAGR and Sharpe use 252-day annualization and Sharpe uses
a zero risk-free rate. Returns include entry fees but no forced final
liquidation. Cash interest, borrowing costs, slippage, and market impact
are excluded.}
\label{tab:final-holdout}
\end{table}

\paragraph{Treynor--Black: weak gains and high turnover.}
Treynor--Black earns 2.70\% of starting capital from trading before paying
6.83\% in fees. At an average daily turnover of 19.20\% of equity,
even the 5 bp charge outweighs its trading gains. Those gains are also
uneven: performance before fees is nearly flat in 2024, positive in 2025,
and negative in 2026 through November 17. Reducing costs would help, but
the signal itself does not produce consistent profits in this sample.
Hedging the estimated factors does not resolve that problem, since a
factor-neutral position can still lose money if residual returns fail to
reverse. The results alone cannot tell us whether estimation noise or a
change in the reversal relationship is responsible.

\paragraph{PCA pairs: lower turnover, modest profitability.}
Pairs earns enough before fees to remain profitable in each of the three
calendar years shown, including the partial year in 2026. Its trading gains
contribute 7.40\% of starting capital, of which 3.37\% is spent on
fees. Average daily turnover is 9.13\%, less than half that of
Treynor--Black. The lower trading burden therefore helps preserve a larger
share of its gains.

The positive result is still small in absolute terms. Pairs uses average
gross exposure of 35.17\%, compared with 100.00\% for Treynor--Black,
because inactive pairs leave part of the portfolio in cash. Borrowing costs
and less favorable execution could reduce the remaining profit further.
For both strategies, the trading gains and fees above are additive
contributions from the simulated portfolio, rather than results from a
separate run without costs.

\begin{figure}[htbp]
\centering
\includegraphics[width=\linewidth]{figures/holdout_performance.png}
\caption{Net cumulative returns for contrarian, factor-hedged Treynor--Black
with a 60-day regression and PCA pairs with a 20-day spread window. Both use
$K=5$ and 5 bp trading fees. Positions start in
cash and begin earning returns after the first holdout close. Missing
observations freeze held positions at their carried marks.}
\label{fig:final-holdout-performance}
\end{figure}

\subsection{Limits of the evaluation}

When a held stock's return is unavailable, its dollar mark stays unchanged
and it cannot trade until observations resume. Final frozen gross exposure
is 14.06\% of equity for Treynor--Black and 1.96\% for pairs.
Unobserved price changes and terminal payouts could alter these returns,
while frozen holdings can weaken portfolio neutrality. This valuation
uncertainty matters especially for Treynor--Black.

\section{Conclusion}

Both PCA strategies are profitable out-of-sample and after transaction costs, but they differ in robustness. Contrarian Treynor--Black shows that estimated alphas reverse rather than persist. It is factor neutral by construction (market beta 0.04) and performs best in crises, but it is concentrated in 2008--09, sensitive to trading costs, and has weakened since 2013. PCA pairs trading is the more robust strategy: it is profitable in 21 of 22 years, well calibrated, insensitive to its parameters, and stronger in the 2018--2023 holdout than before it. Finally, the contrast between in-sample and out-of-sample results shows why rolling estimation matters. Estimating the model on the full sample does not just inflate performance; it reverses the conclusion about which way to trade.

\end{document}
